\chapter*{上编\quad 辑佚}
\phantomsection\addcontentsline{toc}{chapter}{上编\quad 辑佚}
\markboth{上编\quad 辑佚}{上编\quad 辑佚}
\chaptersubtitle{ЧАСТЬ ПЕРВАЯ}{ТЕКСТ}

凡例列于本编之前。所收残简三十七片，分作营建、岁计、撤离、后题四卷，其片数与内容大要略如下表。

\begin{center}
\small
\begin{tabular}{@{}llr@{}}
\toprule
卷次 & 内容 & 片数\\
\midrule
卷一　营建 & 选址、居住环、界、净秽山的前身 & 九\\
卷二　岁计 & 口数、粮秣、水气、工籍、定额 & 十\\
卷三　撤离 & 清点、退役、装载、口令 & 十二\\
卷四　后题 & 后世题记、伪托、讹字、圆片 & 六\\
\bottomrule
\end{tabular}
\end{center}

\vspace{4mm}
\kao{凡\quad 例}

一、释文按原件行款，一字一格。原件无标点，标点为辑校者所加，用以断句；加处一律以文意可通为准。

二、原件所刻为小篆。释文一律改作今字，此为整理之通例，非原件字形；读者若需见原形，须看影印本。书中插图亦同此例，所摹者为释文，非原字。

三、补字以方框（□）标出。凡补字，校记中说明依据；无依据者不补，径作缺字。

四、缺字不可知者作「□」，可推知而字形不全者作「△」。

五、残片编号依基金会编目，自 ФВ-1963-11 起至 ФВ-1963-91 止，号不连续，中间有空号。本书所用简次（简一至简三十七）乃笔者所排定，与编号不同，对照表见附录一。

六、校记只处理字，不处理意思。凡涉及释读之判断，另见下编。

七、凡引传世文献，用中华书局点校本，出注卷次与页码。

八、「本书」指全帙，「残简」指三十七片原件。凡言材料所无，作「残简无一句」；凡言笔者不作解释，作「本书不作解释」。上编分四卷，文中称卷一至卷四，一篇之内互称「卷内」。

九、上编四卷不占章号，故图、表一律用平号（图一、表一），与丛书其余论著通行的「1-1」式不同。二者体例上的差别，源于本书上下两编分编编号，不得已而为之。


\chapter*{卷一\quad 营建}
\phantomsection\addcontentsline{toc}{chapter}{卷一\quad 营建}
\markboth{卷一\quad 营建}{卷一\quad 营建}

卷一九片，无一片为完整文书，形制皆似从账册或规程上裁切而下的断简，前后俱失。九片之中，简一至简六记城池本体，简七至简九记城之边界。此项划分乃笔者所拟定，原件并无任何标记可为佐证；读者若以为简五当归入界一类，亦属可通。

\begin{jian}{简一}
□山之下，去海三舍。地无草木，风自西来，终岁不息。

城中四时温和，无寒暑，昼夜如一。

以石为垣，垣厚四丈，高十二丈。内为环，环各有门，门不可同启。环之外为道，道广十步，环道一周。
\end{jian}
\chuchu{原件残片 11 号，长 21.3 cm，金属，边缘整齐，似为刻后截断；下端有三字蚀尽。}
\xiaojì{「□山」，上字蚀尽，下字作「山」，笔画可辨。此山与简九所称「界外有山」是否一物，二片所记方位不合，详下编第五章。「去海三舍」，三舍之制不可确考；若以三十里为一舍计则九十里，与月面地貌实不相合，「海」字之义附见下编第五章。\par{}「风自西来，终岁不息」，乃城池选址条件之一。月面本无大气流风，此风实出于城内人造气候与定向排风系统；建造期内通风设施已在运转，故当时依实况入录，与月面自然环境无涉，详下编第一章。\par{}「城中四时温和，无寒暑，昼夜如一」十字，为全书论城之立足处：城在月面，月面昼夜温差动辄逾三百度，而环内四时如一，其间必有一套维持系统，然残简对此系统之构造无一言。「昼夜如一」四字，「如」当作「同」解，谓昼与夜全无分别；至于城中究系长明抑或长暗，本片与全书俱无明文，详下编第二章。此十字亦解释了全书记时何以仅用「岁」「月」：既无寒暑可纪，岁便无需分作四时，详下编第一章。\par{}「门不可同启」，谓六门不得同时开启。此属制度规定而非工程限制：同一片内既言「环各有门」，又言「不可同启」，则门之为用不在通行，而在隔绝。}

\begin{jian}{简二}
三环以上，居者稀。一环之口，常不过千。

环尽处为界，界外非所宜往。界内一步，皆在籍。
\end{jian}
\chuchu{原件残片 13 号，长 17.6 cm，陶瓷，正面有釉，背面粗糙；右侧边缘有锯痕。}
\xiaojì{「三环以上」，指居住环三号以上，见简三、简五；城中凡六环，何以三环以上居者稀少，本片未言。「一环之口，常不过千」，此「口」字指门户，非指人口：六环各设一门，门内同时容留者常不足千人，实为通行量之限制。若读「口」为人口，则与简十四「实居十一万九千」相冲，一环不足千则六环全城不足六千，与实居之数相差二十倍，不可通，故本书取前一读，详下编第三章。\par{}「界」字凡四见（简二、简五、简九、简二十一），简九有专述，此处只记其不可逾越。「界内一步，皆在籍」，可见界与籍相应：越界即出籍，出籍即无名。此点与简二十五「无名者不得登」相参，构成下编第八章所论之一大矛盾。}

\begin{jian}{简三}
一环。二环。三环。四环。五环。六环。

环各一门，门各有籍。籍在门吏，不入库。
\end{jian}
\chuchu{原件残片 15 号，长 24.1 cm，金属，字口深，为卷一中刻痕最深一片。}
\xiaojì{六个环名连续刻写，字距均匀，无标点；此种写法见于账簿之目，不见于文书之文，可断此片为某册之目录页。\par{}「籍在门吏，不入库」，为全书所记归档制度之例外。简八言「凡物入库」「板与册不同者以板为正」；此片独言门籍不入库，则门籍别有一套保存成规，至于其法为何，本片未及。}

\begin{jian}{简四}
气自下而上，水自上而下，皆以管。管在壁中，不可见。

凡管，岁一检。检者不入环，于道中为之。检毕，录其数于板，板悬于门。

二十三年，三环西道管裂。检者先见之，环不闭。赏之，以其不泄于众。
\end{jian}
\chuchu{原件残片 18 号，长 19.8 cm，金属，右端断裂，末三字残缺。}
\xiaojì{「检者不入环」，谓检修人员不进入居住环内作业，与简五「外工」之制相合；此类安排在残片中凡三见，并非孤例。「录其数于板，板悬于门」，可与简八「板与册不同者，以板为正」互证：板为当时之物证，悬于门则众目共见，足见气路检修在城中制度中之地位甚重。\par{}下半五行属记事而不属规程，与他片体例异趣：三环西道管裂，检者先发觉之，环未闭而此人受赏，其受赏之由无关管道修葺与否，唯在于「不泄于众」。此为全书唯一一次言及赏赐，所赏者乃一件秘而不宣之事。「二十三年」为何年无从考校，城中纪年与中原异制，本书他处亦无同类纪年可资对校。}

\begin{jian}{简五}
外工三百二十人，居垣外小舍。月一入环，入不过三时。

出则浴，浴水不入渠。衣别贮，不与内同。

外工有越界者，捕得一人。问其故，曰：欲见山。罚如律，籍不除。
\end{jian}
\chuchu{原件残片 21 号，长 22.4 cm，陶瓷，中部有一道横向裂纹，未及字。}
\xiaojì{「外工」之名凡七见，乃城中一类专职工役，居垣外而不与内人杂处；三百二十人之数，与简二十八「外工三百二十人，不入装载之数」相合。「浴水不入渠」谓其洗浴之水另作处置，不与环内水网相接，「衣别贮，不与内同」亦属隔离成规，二句合观，可见城中对「外来之物」防范细密至衣物一级，此与下编第三章所论净秽制度直接相关。\par{}末段记一越界外工被捕，问其故而答曰「欲见山」。此三字为全书三十七片中唯一一句出自城内寻常役夫之口的话，职名与姓氏俱无存，唯留此三字理由。其处置为「罚如律，籍不除」，按律受罚而不削名籍；较诸简十九「逃亡者，籍除，名削」，越界处分反轻，此项差别本书不作解释。}

\begin{jian}{简六}
垣周九里，高十二丈。垣外无墙，以地为限。限外百步，不得置物，不得起舍，不得葬。

限内尝有私舍三区，毁之，其人移籍。
\end{jian}
\chuchu{原件残片 24 号，长 15.2 cm，金属，下端蚀去一角，末二字残。}
\xiaojì{「垣周九里」与简一「垣厚四丈」「高十二丈」当属同一记载之三段，中间脱去一句而无据可补。以九里计，垣周长约当四公里弱，与第 117 全宗第 2187 号卷宗第 12 叶总平面图纸比例尺相合。「不得葬」三字尤为可异：禁令三项中置物、起舍尚属工程之忌，唯独「葬」无工程意义，何以并列其间，本片与他片俱未明言。\par{}末段记限内尝有人私筑三区，既遭拆毁，其人亦随之移籍。于「三区」之外未载他数，亦未及筑者身份；然参以简五「外工居垣外小舍」与简十三「新入者先居外舍三月」，前二者皆属合制之垣外居住，此处则属违禁私建。城中既立严禁而仍有违例建舍之举，此条记载之价值正不在其所言为何，而在揭示制度于实际运作中并未始终贯彻。}

\begin{jian}{简七}
第□环西储物间，架十二，架各五层。层以板隔，板下有记。

凡入此间者，二人俱，一人录，一人取。
\end{jian}
\chuchu{原件残片 28 号，长 11.7 cm，陶瓷。此片即本书三十七片之出土地点之自记。}
\xiaojì{「第□环」，环数字蚀去。此片背面另有刮痕一道，非文字。旧照所见储物间西壁架数与此相合，然层数已不可辨。「二人俱，一人录，一人取」为库房存取之制。三十七片即出自此间，按此制本应另有一册录记，然该册不在所得材料之内，构成全书最大之一处缺环；下编第六章所涉「焚籍」之事，或与此册有关。}

\begin{jian}{简八}
凡物入库，先记于板，后记于册。板与册不同者，以板为正。

板坏则移其字，不移其板。册坏则重录，重录者署其名。
\end{jian}
\chuchu{原件残片 31 号，长 18.9 cm，金属。}
\xiaojì{此条为归档通则，非一次具体记事。板为层板之下所悬小牌，册为总账；二者相歧时以板为正，可见板乃当时原始物证，册为后录之簿。本条与全书编次之原则直接相关，见导论第四节。\par{}「册坏则重录，重录者署其名」为一紧要细节：重录者既须署名，则重录一事即成为可资追查之责任行为。全书三十七片中，唯此一处言及个人之权责，其余皆记制度而不记人。}

\begin{jian}{简九}
界者，非垣非墙，非水非火。在其内者出，必失其半。失者不复。

界外有山，山名不可呼。呼者其罪与越界同。
\end{jian}
\chuchu{原件残片 34 号，长 26.8 cm，金属，为卷一最长一片，下端有一孔，似曾穿绳。}
\xiaojì{「必失其半」语义未明，或指所携之物，或指其人本身；「失者不复」则谓所失不可复得，二句合观，似指涉某种不可逆之减损。\par{}「山名不可呼」与简一所称「□山」相应；此山与下编第五章所论者是否一物，本片未及其用。「呼者其罪与越界同」将言语与越界等量齐观，乃全书唯一将「言语」入罪之条，此点下编第八章另有专论。}
